\documentclass[11pt]{article}
\usepackage[margin=1in]{geometry}
\usepackage{amsmath,amssymb,amsthm}
\usepackage[hidelinks]{hyperref}
\newtheorem{lemma}{Lemma}
\newtheorem{theorem}{Theorem}
\newtheorem{corollary}{Corollary}
\begin{document}
\setlength{\emergencystretch}{1em}
\title{A residue reduction for ordinary support $\{3,7\}$}
\author{}\date{}\maketitle
For every integer $n\ge1$ put
\[
 P_n(t)=\sum_{k=0}^n\binom nk(3n+1)_{4k}t^k.
\]
We prove consecutive coprimality except possibly at
$n\equiv13,45,54\pmod{64}$. Thus the remaining terminal classes
for the ordinary phase $X^3(A+BX^4)$ are
$N\equiv14,46,55\pmod{64}$.
The argument uses an exact mixed-adjoint matrix, a controlled-radius
Frobenius theorem, and a quantitative coefficient consequence of
its nonresonant block. It does not prove noncollision on those
three progressions or an unrestricted Gamma-shape theorem.

All auxiliaries in a fixed pair keep the same parameter
$\gamma=3n+1$:
\[
 B_0=1,\qquad B_{j+1}=(z+4j+\gamma+3)B_j-zB_j'.
\]
Set
\[
 M=7n+1,\qquad
 \kappa=\frac{4(n+1)}{(M+4)(M+5)(M+6)}
        =\frac4{7(7n+5)(7n+6)}.
\]
A cyclic column is the remainder coefficient column modulo
$z^4-u$ in the basis $1,z,z^2,z^3$.
We work over a complete algebraically closed extension $K$ of
$\mathbb Q_2$, with $v(2)=1$, valuation ring $\mathcal O$,
maximal ideal $\mathfrak m$, and residue field $k$.
An overbar means reduction modulo $\mathfrak m$.
In particular a positive valuation is not assumed to be at least one.

\begin{lemma}[The exact actual mixed matrix]\label{lem:mixed-matrix}
For $u\ne0$, the two equations
$P_n(-1/u)=P_{n+1}(-1/u)=0$ are equivalent to rank less than three of
\[
 [B_{n+1},\ B_{n+2}-\kappa uB_n,\
                   B_{n+3}-\kappa(M+11)uB_n](u).        \tag{1}
\]
\end{lemma}
\begin{proof}
Apply \href{DIAGONAL_PRIME_GAP_SFC.tex}{Theorem
\texttt{thm:offset-rank} of the general-offset note}
with $d=4,h=3,\epsilon=-1$.
Its rising-basis polynomial is
\[
 V_3(E)=(E)_2+(M+6)E+(M+5)(M+6).
\]
Thus its three cyclic columns are
$L_0,L_1+\kappa uB_n,L_2+\kappa(M+6)uB_n$.
The exact constant triangularization is
\[
 L_0=B_{n+1},\qquad
 L_1=(M+7)B_{n+1}-B_{n+2},
\]
\[
 L_2=(M+7)(M+6)B_{n+1}-(2M+17)B_{n+2}+B_{n+3}.
\]
Subtract the first-column multiples, change the second column sign,
and eliminate $B_{n+2}$ in the third.
This gives (1) by invertible constant column operations.
The imported equivalence is algebraic in characteristic zero
and includes multiple lower roots; no simplicity assumption is used.
\end{proof}

For later use write
\[
 \mathsf A=3n+4-zD_z+z,\quad
 \Pi_{n,4}(x)=\prod_{\ell=0}^{n-1}(x+4\ell),\quad
 R_i(x)=\prod_{\ell=n}^{n+i-1}(x+4\ell).
\]
Then $B_{n+i}=\Pi_{n,4}(\mathsf A)R_i(\mathsf A)1$.
A dependence in (1) can be normalized by an actual nonzero
coefficient so that $a_1,a_2,a_3\in\mathcal O$ and at least one is a unit.
It gives a zero cyclic column for $\Pi_{n,4}(\mathsf A)Q(\mathsf A)1$,
where
\[
 Q(x)=\sum_{i=1}^3a_iR_i(x)-\kappa uS,\qquad
 S=a_2+(M+11)a_3.                                    \tag{2}
\]
Here $u$ is an evaluated scalar, not a variable acted on by
$\mathsf A$.

\begin{theorem}[Fourteen classes out of sixteen]
\label{thm:fourteen-sixteenths}
If $n\not\equiv6,13\pmod{16}$, then $P_n,P_{n+1}$ are coprime
over $\mathbb Q$ and over $\mathbb C$.
\end{theorem}
\begin{proof}
Every interval of four consecutive positive integers has
two-adic valuation at least three. Hence each coefficient of
$P_n$ at index $k$ has valuation at least $3k$.
The monic reciprocal in $u=-1/t$ implies $v(u)\ge3$
at a root: below that bound the leading term would be uniquely least.

Suppose a common root exists and use (2).
If $v_2(\kappa)>-2$, then $v(\kappa u)>1$.
Since $R_i\bmod\mathfrak m=x^i$,
the polynomial $Q$ is integral and primitive with reduction degree
at most three. Also
\[
 v(R_i(0))=v(4^i(n)_i)\ge2+v_2(n)>1,
\]
so $v(Q(0))>1$.
The exact stage class is
$\mathsf A=c-zD_z+z$, $c=3n+4\in\mathbb Z_2$, seed $1$.
Thus \href{LOWER_TWO_FOURTH_POWER_GAP_SFC.tex}{Theorem
\texttt{thm:boundary} of the fourth-power boundary note},
with $p=2,s=2,q=4,b_s=3,B_s=1$, forbids the zero cyclic column.

Finally
$v_2(\kappa)=2-v_2((7n+5)(7n+6))$.
The condition can fail only when one consecutive factor is
divisible by sixteen, namely at $n\equiv13$ or $6\pmod{16}$.
The polynomials are rational; absence of common roots over $K$
gives a rational Bezout identity, which also holds over $\mathbb C$.
\end{proof}

\begin{corollary}[Fifteen classes out of sixteen]
\label{cor:fifteen-sixteenths}
If $n\not\equiv13,22\pmod{32}$, then $P_n,P_{n+1}$ are coprime.
\end{corollary}
\begin{proof}
Only $n\equiv6,13\pmod{16}$ remain.
For $k\ge2$, counting multiples of $2,4,8$ gives
\[
 v_2((3n+1)_{4k})\ge3k+\lfloor k/2\rfloor>3k.
\]
For $k=1$, an even $n$ contributes $v_2\binom n1\ge1$ in
addition to the baseline three. If $n\equiv13\pmod{16}$,
then $3n+1$ is divisible by eight and its first four-factor
block already has valuation at least four.
Thus every nonleading reciprocal coefficient has valuation
strictly above its value on the line of slope three.
The unique least-term argument excludes $v(u)=3$ as well:
every root has $v(u)>3$ in these two classes.

Put $E=v_2((7n+5)(7n+6))$.
At $E=4$, $v(\kappa)=-2$, whence $v(\kappa u)>1$.
The same boundary proof gives a contradiction.
Only $E\ge5$ remains, equivalent to $n\equiv13,22\pmod{32}$.
\end{proof}

\begin{lemma}[Exact necessity at the second Frobenius stage]
\label{lem:stage-two-necessity}
Use the stage class of the fourth-power boundary theorem:
\[
 A=c-zD_z+F(z),\quad c\in\mathbb Z_2+4\mathcal O,\qquad
 F(z)=z+\sum_{r\ge1}b_rz^{2^r},\quad b_r\in2\mathbb Z_2,
\]
with integral seed $h$ reducing to a constant unit.
Let $n\ge1$, $u\ne0$, $v(u)\ge3$, and let
$Q\in\mathcal O[x]$ be primitive with $\deg\bar Q\le3$.
If the cyclic column of $\Pi_{n,4}(A)Q(A)h$ vanishes, then
\[
 v(Q(0))=1,\qquad \bar Q=a\,x(x-1)^2,\quad a\in k^\times. \tag{3}
\]
The same collision also gives
\[
 v(Q(1))\ge2,\qquad v(Q'(1))\ge1.                       \tag{4}
\]
The ordinary degree of $Q$ need not be bounded.
\end{lemma}
\begin{proof}
Use \href{ARITHMETIC_FRONTIER_TERMINAL_ORDERS.tex}{Lemma
\texttt{lem:artin-hasse-closure} of the arithmetic note}
and the exact block interface in the proof of
\href{LOWER_TWO_FOURTH_POWER_GAP_SFC.tex}{Theorem
\texttt{thm:boundary}}.
There is one nonresonant block $e=1$, of smaller degree two.
After absorption of its unit rising-product factor into the seed,
the first two cyclic vectors are an integral unit basis and all
subsequent vectors are integral.
In the finite divided-Taylor expansion
$\sum_{r\ge0}2^rQ^{[r]}(1)v_r=0$, every term with $r\ge2$
has valuation at least two.
Inversion of that basis gives (4), and hence
\[
 \bar Q=(x-1)^2R,\qquad \deg R\le1.
\]
Its first unit coefficient therefore occurs at index at most one.

Write $Q=\sum c_rx^r$ and choose an actual nonzero coefficient
$\lambda=2^{r_*}c_{r_*}$ of least valuation
\[
 m=\min_r\{v(c_r)+r\}\le1.
\]
The minimum may be fractional. Dividing by this actual field
element makes $Q_*(x)=Q(2x)/\lambda$ integral and primitive.
For $r\ge2$, its coefficient has valuation at least $r-m\ge1$,
so $\deg\bar Q_*\le1$.
The resonant block is exactly a prime-stage collision with this
polynomial, the descended seed, and $v(u_*)\ge1$.

At that prime stage, reduction has $z^2=0$ and a cyclic operator
with eigenvalues $0,1$; the seed has nonzero components in both.
A collision forces $\bar Q_*(1)=0$.
If $m<1$, all its nonconstant coefficients have positive valuation,
so $\bar Q_*$ is a nonzero constant, impossible.
Thus $m=1$, and $\bar Q_*=a(x-1)$.
Its constant coefficient is a unit, giving $v(Q(0))=1$.
Its linear coefficient is a unit as well, so $c_1$ is a unit.
Together with $\bar Q(0)=0$ and the factorization above this gives
$\bar Q=a\,x(x-1)^2$. All normalizations use actual field
elements; no symbol $2^m$ with unspecified fractional exponent is used.
\end{proof}

\begin{lemma}[The actual coefficient constraint]
\label{lem:quantitative-coefficients}
For every common root and every primitive relation (2), the polynomial
$Q$ is integral and
\[
 v(Q(0))=v(\kappa uS)=1,\qquad
 a_1+a_3\in2\mathcal O,\quad a_2\in2\mathcal O.
\]
Consequently, writing $L=v(u)$ and
$E=v_2((7n+5)(7n+6))$,
\[
 \begin{cases}
 L=E-1,&n\text{ odd},\\
 L\le E-2,&n\text{ even}.
 \end{cases}                                             \tag{5}
\]
\end{lemma}
\begin{proof}
All nonconstant coefficients of $Q$ are integral, and some are
units, because $\bar R_i=x^i$ and at least one $a_i$ is a unit.
Only its constant coefficient could be nonintegral.
If $v(Q(0))<0$, divide by the actual field element $Q(0)$.
The resulting polynomial is integral and primitive with reduction
the nonzero constant one. Its zero cyclic column contradicts
\href{ARITHMETIC_FRONTIER_TERMINAL_ORDERS.tex}{Lemma
\texttt{lem:frobenius-noncollision}}, since its reduction degree
is zero, less than three.
Thus $Q$ is integral and primitive, and
Lemma~\ref{lem:stage-two-necessity} applies.

The values $R_i(0)$ have valuation at least two.
Equation (2) and $v(Q(0))=1$ therefore imply
$v(\kappa uS)=1$.
The integer polynomials also satisfy
\[
 R_i(1)\equiv1\pmod4,\quad R_1'(1)=1,\quad
 R_2'(1)\equiv0\pmod2,\quad R_3'(1)\equiv1\pmod2.
\]
The derivative estimate (4) implies $a_1+a_3\in2\mathcal O$.
The value estimate gives
$a_1+a_2+a_3-\kappa uS\in4\mathcal O$.
As $\kappa uS\in2\mathcal O$, it follows that $a_2\in2\mathcal O$.
This conclusion uses quantitative estimates; reduction modulo
$\mathfrak m$ alone would not imply divisibility by two.

By (3), $a_3$ is a unit.
If $n$ is odd, $M+11=7n+12$ is odd, so $S$ is a unit.
If $n$ is even, $S\in2\mathcal O$, so $v(S)\ge1$.
Since $1=v(\kappa uS)=2-E+L+v(S)$, both cases in (5) follow.
\end{proof}

\begin{corollary}[Sixty-one terminal classes out of sixty-four]
\label{cor:sixty-one-sixty-four}
For every $n\ge1$ with
$n\not\equiv13,45,54\pmod{64}$, the polynomials
$P_n,P_{n+1}$ are coprime over $\mathbb Q$ and over $\mathbb C$.
Consequently every nonzero ordinary phase $X^3(A+BX^4)$ has
a nonzero factorial moment in each consecutive positive pair
with terminal $N\not\equiv14,46,55\pmod{64}$.
\end{corollary}
\begin{proof}
By Corollary~\ref{cor:fifteen-sixteenths}, odd remaining indices
satisfy $n\equiv13\pmod{32}$.
For an even remaining index, $E\ge5$ and the strict-radius
argument above gives $L>3$.
If $E=5$, (5) gives $L\le3$, a contradiction.
Hence an even remaining index must have $E\ge6$,
which is exactly $n\equiv54\pmod{64}$.
This proves the stated three-class residual set.

For $AB\ne0$ the order-$j$ factorial moment is
$A^j(3j)!P_j(B/A)$; put $N=n+1$.
If one coefficient vanishes, the nonzero phase is a monomial
and every positive factorial moment is nonzero.
\end{proof}

\paragraph{Evidence and limits.}
The exact checker verifies complete generic column identities,
the full periods16,32,64, generic value/derivative congruences,
the valuation bounds and literal cyclic-column controls.
The six written proofs supply the unbounded quantifiers.
The three remaining progressions and unrestricted binomial SFC
are open. Internal mathematical audit and replay do not constitute
independent whole-proof replay, formal verification, external review
or a priority claim.
\end{document}
